\documentclass{article}

% Language setting
% Replace `english' with e.g. `spanish' to change the document language
\usepackage[english]{babel}

% Set page size and margins
% Replace `letterpaper' with `a4paper' for UK/EU standard size
\usepackage[letterpaper,top=2cm,bottom=2cm,left=3cm,right=3cm,marginparwidth=1.75cm]{geometry}

% Useful packages
\usepackage{amsmath}
\usepackage{graphicx}
\usepackage[colorlinks=true, allcolors=blue]{hyperref}
\usepackage{lipsum} 


\newcommand{\titreproposefr}{Le titre}
\newcommand{\titreproposeen}{The title}
\newcommand{\themefr}{Le thème en une phrase.}
\newcommand{\themeen}{The theme in one sentence.}

\author{You}
%\newcommand{\coredac}{mon ou ma co-redacteur·rice}



\title{Soumission de projet de numéro thématique}


\begin{document}
\maketitle



\section{Présentation de la thématique}

\begin{tabular}{|rp{10cm}|}
\hline
\textbf{Titre} : & \titreproposefr \\
\textbf{Title} : & \titreproposeen\\
\hline
\textbf{Thème} :& \themefr\\
\textbf{Theme} :& \themeen\\
\hline
\textbf{Orientation} : & choisir : thématique ou vers des applications  \\
\hline
\textbf{Mots clés} : & 3-5 mots clés \\
\textbf{Keywords} : & 3-5 keywords \\
\hline

\end{tabular}
\\



\subsection{Proposition d'appel à contribution}

% Écrire la partie de l'appel à contribution présentant le thème. Une description des enjeux au niveau de l'état de l'art et une description précise du thème sont attendues.
% La présentation de la revue et les modalités seront ajoutées après acceptation de la proposition.

\lipsum[1-2]


\subsection{Activités récentes sur ce thème}
% Une présentation succincte des activités passées récentes sur ce thème (journées d’étude, workshop national ou international, etc.), si possible avec le nombre de contributions.

\lipsum[1]

\section{Équipe pour la coordination du numéro}

\subsection{Coordination du numéro}
% Notez que nous serons attentifs aux équilibres Femme-Homme 

% Présenter les co-rédacteurs ou co-rédactrice en cheffe qui pourraient participer au numéro
% Préciser leurs coordonnées, affiliation, thèmes de recherche, et faits marquants. Expliquer en un court paragraphe pourquoi cela ferait sens que vous trailliez ensemble.

\begin{itemize}
\item Mme X, université A
\item M. Y, université B
\end{itemize}



\subsection{Comité de rédaction extérieur}
% Notez que nous serons attentifs aux équilibres Femme-Homme 

%Donner sous forme de liste les noms, affiliation et sujet de recherche des potentiels membres de votre comité de programme externe.

\begin{itemize}
\item M Z, université C
\item Mme U, université D
\end{itemize}

\nocite*
\bibliographystyle{alpha}
\bibliography{sample}

\end{document}